\scanpage{113}
$\tau_\omega(H^{**}(\mathscr X))$ [qui est une algèbre graduée \deleted{\unsure{asso}} alternée, entre autres, de $M_\omega$] est \deleted{provient d'un} dans l'image essentielle de $\mathrm{Mod}(k)$, i.e. peut s'écrire (de façon ess. unique \unsure{à} $\unsure{\cong}$ \unsure{au niveau module|ou encore usuelle})
\[
 \mathscr H_C^{**}(\mathscr X)=\deleted{H}\!H_C^*(X)\otimes_k S_k
\]
[ et on sait que $H_C^*(X)$ hérite alors la structure de $\unsure{k'|k|l'}$-algèbre graduée de $\tau_\omega(H^{**}(\mathscr X))$. Cela implique que \deleted{pour} les $\deleted{\unsure{A}}C_n^{**}$ ont la $=$ propriété (faire $\mathscr X=\Delta_n$).
\deleted{On veut de plus que \inserted{$\forall\mathscr X$,} l'augmentation $S_k\to C^{**}(\mathscr X)$ définisse} On va préciser cette exigence

On peut considérer $\deleted{\unsure{A}}C_n^{**}$ comme un faisceau (à structure de type $\mathscr D$) sur \deleted{\unsure{la}} le topos $(\mathrm{Ss})$, il est coaugmenté par le faisceau constant $S_k$
\[
 S_k\longrightarrow C_{n}^{**}
\]
\deleted{\unsure{d'où}} [NB, on regarde ici le degré extérieur comme degré principal] de sorte qu'on a un diagramme (dans \deleted{\unsure{la}} cat. dérivée \deleted{$D(S_k$} \deleted{\unsure{consistant d'objets}} \inserted{$D(S_k\text{-mod gradué})$}) :
\diagram{\begin{tikzcd}[ampersand replacement=\&,column sep=small]
 \mathbb R\Gamma^{\unsure{*}}(\mathscr X,S_k)\arrow[r] \& \mathbb R^{*}\Gamma(\mathscr X,C_{\unsure{*|n}}^{**})\\
 |[alias=tag]|{\text{(1)}} \& C^{**}(\mathscr X)=\Gamma(\mathscr X,C^{**}\arrow[u]
\end{tikzcd}}
% archive sheet number: 112
